\section{Chicago and the Evidence of Local Variation}

Chicago became one of the Institute's most visible American centers after Cardinal Francis George invited it in 2004 to restore a church then known as Saint Gelasius and develop the Shrine of Christ the King. A serious fire in 2015 interrupted restoration but did not end the project. On 28 February 2016 the Archdiocese of Chicago announced that it had deeded the shrine building and land to the Institute, which committed to stabilize and restore them. This official record fixes the property setting without disclosing every agreement governing the apostolate. After \emph{Traditionis custodes}, the shrine became evidence of a changed local relationship.

Chicago's archdiocesan newspaper published Cardinal Blase Cupich's general implementation policy on 5 January 2022, effective 25 January. It required the archbishop's permission for private or public celebration with the 1962 Missal and written affirmations concerning the liturgical reform, Vatican II, and the papal magisterium. The policy supplies controlling local context but does not name the shrine or itself explain the later event. The current ICKSP Chicago page states one precise fact: celebration of public Masses has been suspended since 1 August 2022. Contemporary Catholic News Agency reporting described conflicting explanations and uncertainty about the parties' positions. The checked set contains no public individual decree, full correspondence, or executed agreement resolving cause and scope.

The checked record therefore does not establish that the Institute was expelled from Chicago, globally banned, canonically suppressed, or proved unwilling to obey. Nor does it support minimizing the suspension: public Mass was the central visible work of the shrine. A source-first account states the material loss of public celebration and preserves the missing file.

\statusframe{Archdiocese of Chicago property release; Chicago's general 2022 liturgical policy; ICKSP Chicago page, current through 17 July 2026; Catholic News Agency report.}{The archdiocese announced its 2016 deed of the building and land to the Institute; archdiocesan permission governs use of the 1962 Missal; public Masses at the shrine have been suspended since 1 August 2022 amid publicly reported disagreement.}{Neither the property record nor the general policy supplies the missing individual act. No checked decree or correspondence establishes current title, expulsion, complete cessation of every activity, a definitive motive, or an adjudicated breach.}

\subsection{Later reception in other dioceses}

Subsequent official diocesan records show that Chicago did not establish a universal rule against ICKSP ministry. On 14 September 2023 the Diocese of Worcester published a decree establishing Saint Paul Oratory at Warren in a local reorganization; contemporary diocesan reporting and the continuing work identify ICKSP service there. In spring 2025 the Diocese of Syracuse announced that the Institute would send a priest and brother to serve Saint Mary of the Assumption Shrine in Oswego beginning in June. A July 2026 diocesan statement still identified the shrine as diocesan and directed by the Institute.

In July 2025 the Bishop of Saint Petersburg announced, on presentation by the ICKSP United States provincial, the appointment of an ICKSP canon as rector of Epiphany of Our Lord Shrine in Tampa effective 1 September. The current diocesan directory preserves the location. These are affirmative local acts after \emph{Traditionis custodes}; they do not create a universal exemption any more than Chicago creates a universal prohibition.

\subsection{The constitutional ecology made visible}

The contrast demonstrates a canonical structure reported at the 1990 diocesan origin and explicit in universal law and the 2008 act. A pontifical-right society supplies members and internal government. A bishop governs public worship and apostolate in his diocese and receives, restricts, or structures a work under applicable law. Roman law shapes all of them, but application reaches the faithful through particular acts and agreements.

Local variation can be pastorally painful and legally contested without proving incoherence in the category itself. It is the expected place where universal legislation, proper law, episcopal judgment, personnel, property, and local need meet. The history must therefore remain location-specific.

\claimcheck{``Chicago shows that the ICKSP was suppressed after \emph{Traditionis custodes}.''}{The Institute continued internationally and was later received in other United States dioceses. Chicago proves a local suspension of public Masses, not dissolution of the society or a universal ban.}
